\section{Legal Postures}

Search can support posture claims such as transparency, robustness,
least-change comparison, alternative-plan availability, or auditability. Each
posture requires different evidence and should be tied to specific commands,
artifacts, and limitations.

\begin{table}[htbp]
\centering
\small
\begin{tabular}{p{0.22\linewidth}p{0.34\linewidth}p{0.28\linewidth}}
\toprule
Posture & Supporting record & Common overclaim \\
\midrule
Transparency & Parameters, seeds, constraints, and RCTX & The process is neutral \\
Robustness & Alternative plans under declared protocols & The chosen plan is legally safe \\
Least change & Baseline, distance metric, and comparison set & The change is justified by metric alone \\
Alternative availability & Generated candidate set and filters & Any alternative is equally viable \\
Auditability & RPLAN/RCTX, certificates, and manifests & Audit pass proves compliance \\
\bottomrule
\end{tabular}
\caption{Legal/policy postures and the records they require.}
\end{table}

The posture vocabulary is useful because it keeps computational artifacts in
their lane. A search workflow can show that many alternatives were considered,
or that a selected plan survived a declared audit profile. It cannot decide
which legal standard applies, which evidence a court or agency will credit, or
whether unmodeled facts change the conclusion.
